% ============================================================
\chapter{Fenchel--Moreau--Rockafellar Duality}
\label{ch:fenchel}
% ============================================================

The Legendre transform, known to physicists since the eighteenth century for passing from the Lagrangian to the Hamiltonian, was generalised by Werner Fenchel in the 1940s into a tool of remarkable power: \emph{convex conjugation}. The idea is to represent a convex function not by its graph, but by the collection of its supporting hyperplanes---a kind of geometric duality where every convex function has an alter ego. The Fenchel-Moreau biconjugation theorem states that for a lower semicontinuous convex function, applying this transformation twice returns the original function. This involution is the key to duality in convex optimisation.

\section{Introduction}

Convex conjugation, or the Legendre--Fenchel transform, is a fundamental tool of
convex analysis. It establishes a correspondence between convex functions and enables
the construction of powerful dual problems. The Fenchel--Moreau biconjugation theorem
characterizes lower semicontinuous convex functions.

\section{Fenchel Conjugate}

\begin{definition}[Fenchel Conjugate]
Let $f : \R^n \to \R \cup \{+\infty\}$. The \emph{Fenchel conjugate} (or
\emph{Legendre--Fenchel transform}) of $f$ is:
\[
    f^*(y) = \sup_{x \in \R^n} \left\{\inner{y}{x} - f(x)\right\}.
\]
\end{definition}

\begin{proposition}[Immediate Properties]
\begin{enumerate}
    \item $f^*$ is always convex and lower semicontinuous (as a supremum of affine
          functions in $y$).
    \item \textbf{Fenchel--Young inequality}:
    $f(x) + f^*(y) \geq \inner{x}{y}$ for all $x, y$.
    \item If $f$ is proper, then $f^*$ never takes the value $-\infty$.
\end{enumerate}
\end{proposition}

\begin{proof}[Proof of the Fenchel--Young inequality]
By definition of $f^*$:
\[
    f^*(y) = \sup_{z} \{\inner{y}{z} - f(z)\} \geq \inner{y}{x} - f(x),
\]
hence $f(x) + f^*(y) \geq \inner{x}{y}$.
\end{proof}

\begin{keyformulas}[title=\textbf{Fenchel--Young Inequality}]
\[
    f(x) + f^*(y) \geq \inner{x}{y}, \quad \forall x, y \in \R^n,
\]
with equality if and only if $y \in \partial f(x)$ (or equivalently
$x \in \partial f^*(y)$).
\end{keyformulas}

\section{Examples of Conjugates}

\begin{example}[Classical Conjugates]
\begin{enumerate}
    \item \textbf{Affine function}: $f(x) = \inner{a}{x} + b$.
    Then $f^*(y) = \begin{cases} -b & \text{if } y = a, \\ +\infty & \text{otherwise.}\end{cases}$

    \item \textbf{Quadratic}: $f(x) = \frac{1}{2}x^TQx$ with $Q \succ 0$.
    Then $f^*(y) = \frac{1}{2}y^TQ^{-1}y$.

    \item \textbf{Norm}: $f(x) = \norm{x}_p$ for $p \geq 1$.
    Then $f^* = \iota_{B_q}$ where $B_q = \{y \mid \norm{y}_q \leq 1\}$ and
    $1/p + 1/q = 1$.

    \item \textbf{Squared norm}: $f(x) = \frac{1}{2}\norm{x}^2$.
    Then $f^*(y) = \frac{1}{2}\norm{y}^2$ (self-conjugate).

    \item \textbf{Indicator}: $f = \iota_C$.
    Then $f^* = \sigma_C$ (support function).

    \item \textbf{Negative entropy}: $f(x) = x\log x - x$ for $x > 0$.
    Then $f^*(y) = e^y$.

    \item \textbf{Log-sum-exp}: $f(x) = \log(\sum_i e^{x_i})$.
    Then $f^*(y) = \sum_i y_i \log y_i$ if $y \in \Delta_n$, $+\infty$ otherwise.
\end{enumerate}
\end{example}

\begin{proof}[Proof for the quadratic]
$f^*(y) = \sup_x \{\inner{y}{x} - \frac{1}{2}x^TQx\}$. The optimality condition
gives $y - Qx = 0$, i.e., $x = Q^{-1}y$. Thus:
\[
    f^*(y) = \inner{y}{Q^{-1}y} - \frac{1}{2}(Q^{-1}y)^TQ(Q^{-1}y)
    = y^TQ^{-1}y - \frac{1}{2}y^TQ^{-1}y = \frac{1}{2}y^TQ^{-1}y.
\]
\end{proof}

\section{Fenchel--Moreau Biconjugation Theorem}

\begin{theorem}[Fenchel--Moreau]
\label{thm:fenchel_moreau}
Let $f : \R^n \to \R \cup \{+\infty\}$ be proper. Then:
\[
    f^{**} = f \quad \Longleftrightarrow \quad
    f \text{ is convex and lower semicontinuous.}
\]
In general, $f^{**}$ is the greatest convex l.s.c.\ function majorized by $f$
(the \emph{convex l.s.c.\ regularization} of $f$, denoted $\overline{\mathrm{conv}}(f)$).
\end{theorem}

\begin{proof}[Proof sketch]
First show that $f^{**} \leq f$ (from the Fenchel--Young inequality). Then use the
separation theorem: if $f$ is convex l.s.c.\ and $(x_0, t_0) \notin \mathrm{epi}(f)$,
there exists a hyperplane separating $(x_0, t_0)$ from $\mathrm{epi}(f)$, which shows
$f^{**}(x_0) \geq t_0$. Since this holds for all $t_0 < f(x_0)$, we get $f^{**} \geq f$.
\end{proof}

\section{Link Between Conjugate and Subdifferential}

\begin{theorem}[Conjugate--Subdifferential Equivalence]
\label{thm:conjugate_subdiff}
Let $f$ be proper convex and l.s.c. The following are equivalent:
\begin{enumerate}
    \item $y \in \partial f(x)$,
    \item $x \in \partial f^*(y)$,
    \item $f(x) + f^*(y) = \inner{x}{y}$.
\end{enumerate}
\end{theorem}

\begin{proof}
$(1) \Rightarrow (3)$: If $y \in \partial f(x)$, then for all $z$:
$f(z) \geq f(x) + \inner{y}{z-x}$, i.e., $\inner{y}{z} - f(z) \leq \inner{y}{x} - f(x)$.
So $f^*(y) = \inner{y}{x} - f(x)$.

$(3) \Rightarrow (1)$: $f^*(y) = \inner{y}{x} - f(x)$ means the supremum in the
definition of $f^*$ is attained at $x$, so $y \in \partial f(x)$.

$(1) \Leftrightarrow (2)$: By $f^{**} = f$ and symmetry.
\end{proof}

\section{Fenchel--Rockafellar Duality}

\begin{theorem}[Fenchel--Rockafellar Duality]
\label{thm:fenchel_rockafellar}
Let $f, g : \R^n \to \R \cup \{+\infty\}$ be proper convex and $A \in \R^{m \times n}$.
Consider the primal problem:
\[
    (P) \quad p^* = \inf_{x \in \R^n} \{f(x) + g(Ax)\}
\]
and the dual problem:
\[
    (D) \quad d^* = \sup_{y \in \R^m} \{-f^*(-A^Ty) - g^*(y)\}.
\]
Then:
\begin{enumerate}
    \item (Weak duality) $d^* \leq p^*$.
    \item (Strong duality) If $\mathrm{ri}(\mathrm{dom}(f)) \cap
          A^{-1}(\mathrm{ri}(\mathrm{dom}(g))) \neq \emptyset$ (constraint
          qualification), then $d^* = p^*$ and the dual supremum is attained.
\end{enumerate}
\end{theorem}

\begin{proof}[Proof of weak duality]
For all $x$ and $y$, by the Fenchel--Young inequality:
\begin{align*}
    f(x) &\geq \inner{-A^Ty}{x} - f^*(-A^Ty) = -\inner{y}{Ax} - f^*(-A^Ty), \\
    g(Ax) &\geq \inner{y}{Ax} - g^*(y).
\end{align*}
Summing: $f(x) + g(Ax) \geq -f^*(-A^Ty) - g^*(y)$.
Taking the infimum over $x$ and supremum over $y$: $p^* \geq d^*$.
\end{proof}

\section{Special Case: $A = I$}

\begin{corollary}[Fenchel Duality]
For proper convex $f, g$ with $\mathrm{ri}(\mathrm{dom}(f)) \cap
\mathrm{ri}(\mathrm{dom}(g)) \neq \emptyset$:
\[
    \inf_x \{f(x) + g(x)\} = \sup_y \{-f^*(-y) - g^*(y)\}
    = -\inf_y \{f^*(-y) + g^*(y)\}.
\]
\end{corollary}

\section{Conjugate Calculus Rules}

\begin{proposition}[Calculus Rules]
\begin{enumerate}
    \item \textbf{Separability}: if $f(x) = \sum_i f_i(x_i)$, then
          $f^*(y) = \sum_i f_i^*(y_i)$.
    \item \textbf{Scaling}: $(\alpha f)^*(y) = \alpha f^*(y/\alpha)$ for $\alpha > 0$.
    \item \textbf{Translation}: if $g(x) = f(x - a)$, then $g^*(y) = f^*(y) + \inner{y}{a}$.
    \item \textbf{Adding a linear term}: if $g(x) = f(x) + \inner{b}{x} + c$,
          then $g^*(y) = f^*(y - b) - c$.
    \item \textbf{Linear composition}: if $g(x) = f(Ax)$ with $A$ invertible,
          then $g^*(y) = f^*(A^{-T}y)$.
    \item \textbf{Infimal convolution}: $(f \mathop{\square} g)^* = f^* + g^*$.
\end{enumerate}
\end{proposition}

\section{Application: Dual Formulation of LASSO}

The primal LASSO problem is:
\[
    \min_x \frac{1}{2}\norm{Ax - b}^2 + \lambda\norm{x}_1.
\]

Writing $f(x) = \lambda\norm{x}_1$ and $g(z) = \frac{1}{2}\norm{z - b}^2$ with $z = Ax$:

\begin{align*}
    f^*(y) &= \iota_{\{y \mid \norm{y}_\infty \leq \lambda\}}(y), \\
    g^*(w) &= \frac{1}{2}\norm{w}^2 + \inner{w}{b}.
\end{align*}

Fenchel--Rockafellar duality gives:
\[
    \max_\nu \left\{-\frac{1}{2}\norm{b - \nu}^2\right\}
    \quad \text{subject to} \quad \norm{A^T\nu}_\infty \leq \lambda.
\]

\section{Python Implementation}

\begin{codeblock}[title=\textbf{Computing conjugates and verifying Fenchel--Young}]
\begin{minted}{python}
import numpy as np
import cvxpy as cp

def fenchel_conjugate_numerical(f, y, n):
    """Numerical computation of f*(y) = sup_x {<y,x> - f(x)}."""
    x = cp.Variable(n)
    objective = cp.Maximize(y @ x - f(x))
    prob = cp.Problem(objective)
    prob.solve()
    return prob.value

# Example: f(x) = ||x||^2 / 2
n = 5
y_test = np.random.randn(n)

f = lambda x: 0.5 * cp.sum_squares(x)
f_star_val = fenchel_conjugate_numerical(f, y_test, n)
f_star_exact = 0.5 * np.sum(y_test**2)

print(f"f*(y) numerical:  {f_star_val:.6f}")
print(f"f*(y) analytical: {f_star_exact:.6f}")

# Verify Fenchel-Young: f(x) + f*(y) >= <x,y>
x_test = np.random.randn(n)
f_x = 0.5 * np.sum(x_test**2)
lhs = f_x + f_star_exact
rhs = np.dot(x_test, y_test)
print(f"f(x)+f*(y) = {lhs:.4f} >= <x,y> = {rhs:.4f}: {lhs >= rhs-1e-10}")

# Fenchel duality for LASSO
m, nn = 50, 20
A = np.random.randn(m, nn)
b = np.random.randn(m)
lam = 1.0

# Primal
x = cp.Variable(nn)
primal = cp.Problem(cp.Minimize(
    0.5 * cp.sum_squares(A @ x - b) + lam * cp.norm1(x)))
primal.solve()
print(f"\nLASSO primal: {primal.value:.6f}")

# Dual
nu = cp.Variable(m)
dual = cp.Problem(cp.Maximize(
    -0.5 * cp.sum_squares(b - nu)),
    [cp.norm_inf(A.T @ nu) <= lam])
dual.solve()
print(f"LASSO dual:   {dual.value + 0.5*np.sum(b**2):.6f}")
\end{minted}
\end{codeblock}

\section{Exercises}

\begin{exercise}[$\star$]
Compute the conjugate of $f(x) = \frac{1}{p}\norm{x}_p^p$ for $p > 1$. Verify that
$f^*(y) = \frac{1}{q}\norm{y}_q^q$ with $1/p + 1/q = 1$.
\end{exercise}

\begin{exercise}[$\star$]
Show that if $f(x) = \iota_C(x)$, then $f^* = \sigma_C$.
\end{exercise}

\begin{exercise}[$\star\star$]
Show that $f^{**} \leq f$ always, and that $f^{**} = f$ if and only if $f$ is convex
and lower semicontinuous.
\end{exercise}

\begin{exercise}[$\star\star$]
Compute the conjugate of $f(x) = \log(1 + e^x)$ and interpret the result in terms
of entropy.
\end{exercise}

\begin{exercise}[$\star\star$]
Derive the dual problem of the LASSO via Fenchel--Rockafellar duality.
\end{exercise}

\begin{exercise}[$\star\star\star$]
Prove the Fenchel--Moreau theorem using the hyperplane separation theorem.
\end{exercise}

\begin{exercise}[$\star\star\star$]
Show the rule $(f \mathop{\square} g)^* = f^* + g^*$ and deduce that the infimal
convolution of two proper convex l.s.c.\ functions is convex l.s.c.
\end{exercise}
